% ch14 -- Open questionn und Forschungslinien
\label{ch:opene_fragen}
% Source documente: A230, 345

\epigraph{A theory that has no open questions has stopped doing research.}{}

\section{Open Questions in the Hyperbit--FFGFT Connection}

\textbf{Cl(6) und GF(27).}
Is the Clifford algebra $\mathrm{Cl}(6)$ a special case of the FFGFT Galois structure
auf $T^4$, or are both independent paths to the same physics? Falls es eine
Abbildung gibt: Wie lauten die Übersetzungsregeln?

\textbf{Universal bit mass.}
Matzkes $\mbit = \mP \ln 2 / (2\pi)$ ist eine universelle Konstante ---
in der FFGFT ist sie scale-sensitive. Is there a regime, in dem beide
Vorhersagen experimentally distinguishable sind?

\textbf{p-Bit-Vorhersage.}
FFGFT predicts: natural p-bit materials bei $L_8 \approx 22\,\text{nm}$.
Gibt es Messungen, die diese Skala als besonders stabil oder instabil zeigen?

\section{Open Questions in the Landauer Interpretation}

\textbf{multiple-realisation invariance (Dok.\,A272).}
Gibt es physikalische Systeme, in denen die Invarianz bricht --- Systeme, in
denen die Energie mit dem Informationsgehalt mitvariiert durch konstruktive
Kopplung? Das would affect the proof of the carrier theorem, aber nicht
seine Prämisse.

\textbf{Phasen-Dissipation.}
Is the information loss in the phase in a quantum measurement
thermodynamically equivalent to the magnitude information loss? (Dok.\,290)

\textbf{Kontinuierlicher Entropiestrom.}
DRAM-Refresh-Kosten: How precisely does the Landauer bound agree mit gemessenen
Entropieströmen in modernen Prozessoren überein? Gibt es systematische
Abweichungen durch Quantenkorrekturen?

\section{Open Question: Matzke's Axiom}

\textbf{Warum $e_i^2 = +1$?}
In Matzke's Hyperbit framework it remains open, warum genau dieses Vorzeichen
die Basis von Physik bildet. In der FFGFT ist die analoge Frage beantwortet:
$\xi = 4/30000$ folgt algebraically aus den Galois group orders (Dok.\,338).
Matzkes opene Frage has no direct counterpart in the FFGFT.

\section{Experimental Approaches}

Three experimental lines are particularly promising:
\begin{enumerate}
  \item \textbf{Bérut-Typ-Experimente} auf p-Bit-Skala: Kann man Landauers Grenze
    bei $L_8 \approx 22\,\text{nm}$ and room temperature be measured directly?
  \item \textbf{Qubit-Dekohärenz-Kalorimetrie}: Does the heat generation during
    decoherence agree with the Landauer formula?
  \item \textbf{Galois-Spektroskopie}: Does $\mathrm{GF}(27)$-Struktur in
    Teilchenspektren (Dok.\,356--357) show measurable signatures?
\end{enumerate}

\section{Status Balance}

\begin{center}
\small
\begin{longtable}{@{}p{0.72\textwidth}l@{}}
\toprule
\textbf{Statement} & \textbf{Status} \\
\midrule
\endfirsthead
\toprule
\textbf{Statement} & \textbf{Status} \\
\midrule
\endhead
Bit $\neq$ physical region (Shannon vs.\ Boltzmann) & \status{B} \\
$S = k_B \ln W$ as starting point & \status{SETZUNG} \\
Region reduction $\Rightarrow$ Landauer bound & \status{B} \\
Liouville-Theorem & \status{B} \\
Content neutrality of thermodynamics & \status{B} \\
$h^f$ as fundamental discretisation & \status{B} \\
Hardware $\neq$ Rechenoperation & \status{B} \\
Energy is a property of the carrier, not the information & \status{B} \\
Permanentspeicher as counter-experiment to Vopson & \status{B} \\
$\Ebit = \hbar c / L$ (scale-sensitive) & \status{B} \\
$\xi = 4/30000$ from Galois group orders & \status{K} \\
Matzkes $\mbit = \mP \ln 2/(2\pi)$ (Landauer bei $T_P$) & \status{SETZUNG} \\
Topological information conservation in FFGFT & \status{K} \\
$n_{\text{thresh}} = 6{,}41$ system-dependent (not universal) & \status{K} \\
$T^4$-Kompaktifizierung resolves the Hawking paradox & \status{B} \\
Matzke's question $e_i^2 = +1$: open & \status{H} \\
multiple-realisation invariance: open & \status{H} \\
Bérut-type experiment at p-bit scale & \status{S} \\
\bottomrule
\end{longtable}
\end{center}
